\newpage
\section{Performance Analysis}
\label{sec:baseline_analysis}

% Chapter purpose: Prove the problem exists.
% Present empirical evidence that always-on module activation is measurably inefficient.
% Three dimensions: energy (5.1), latency (5.2), quality (5.3), then synthesis (5.4).
% Figures source: src/lkc-llm/scripts/experiments/hotpot/2-performance/output/plots/
%                 src/lkc-llm/scripts/experiments/hotpot/3-quality/output/plots/


\subsection{Energy Breakdown}
\label{subsec:energy-breakdown}

% --- Experimental anchor ---

Earlier, we introduced the hypothesis that the optional RAG stages
require more energy, implying that not all of the options need to be active for
every query to produce acceptable answers. To test this claim, we executed all
1{'}200~queries under each of the eight configurations, producing
9{'}600~sessions under the protocol defined earlier.
This section presents the energy dimension of
that experiment. The main comparison is between C0 (minimal pipeline) and C7
(all options enabled); single-module configurations (C1, C2, C3) isolate
per-modules contributions. Table~\ref{tab:configs} lists all eight
configurations and their module settings.

% --- Total energy per request and Energy/token

The figure below reports the mean total energy per request across
all eight configurations. The baseline consumes 318~J per request on average, while the
KMS pipeline consumes 1{'}786~J---a ratio of 5.6$\times$.

Energy is measured as counter deltas over session windows without subtracting
idle baseline power. During HyDE's additional 4.3\,s, the embedding and
database nodes accumulate an estimated $\sim$70\,J of idle draw---approximately
5\% of HyDE's 1{'}355\,J marginal cost. The qualitative finding is unchanged.
\begin{figure}[h!]
    \centering
    \includegraphics[width=0.85\textwidth]{total_energy_by_config.png}
    \caption{Mean total energy per request across configurations}
    \label{fig:total-energy}
    % Source: hotpot/2-performance/output/plots/total_energy_by_config.png
\end{figure}

Two energy clusters emerge:
\begin{itemize}
    \item Configurations that include HyDE (C1, C4, C5, C7) form a high-energy group 
    between 1{'}666--1{'}786~J.
    \item Configurations not including HyDE which cluster below 435~J, with hybrid retrieval (C3)
    nearly identical to C0 (324~J against 318~J), adding only 2\% energy overhead.
\end{itemize}

Module costs are approximately additive given the hybrid search influence.
C7 (1{'}786~J) exceeds C4 (1{'}783~J) by only 3~J, closely matching the 6~J cost of hybrid retrieval
alone (C3$-$C0). Given that, the per-module analysis in the next section characterises
each module's energy overhead, focusing on the cross-module comparisons.

% --- Stage-level energy attribution ---

Table~\ref{tab:stage-energy} details the five measured stages for C0 and C7.
The prompt augmentation stage (HyDE) is the largest single contributor in the
full pipeline, consuming 1{'}271~J---4.0$\times$ the entire baseline (C0)
request. Reranking adds 111~J. Answer generation remains roughly constant
across configurations (268--302~J), indicating that the LLM generation phase
does not vary with retrieval enrichment.

\begin{table}[h!]
\centering
\begin{tabular}{l r r}
\hline
\textbf{Pipeline stage} & \textbf{Baseline (J)} & \textbf{KMS Pipeline (J)} \\
\hline
Embedding           &  43.2 &   63.6 \\
HyDE (prompt aug.)  &   --- & 1354.5 \\
Retrieval (DB)      &   7.0 &   38.1 \\
Reranking           &   --- &  105.8 \\
Generation          & 267.7 &  302.4 \\
\hline
\end{tabular}
\caption{Stage-level energy attribution (C0 vs C7).}
\label{tab:stage-energy}
\end{table}

% --- Feature attribution ---

The differential view $\Delta E = E(\text{C}_i) - E(\text{C0})$ quantifies
each module's mean marginal cost (Figure~\ref{fig:feature-attribution}). HyDE adds
1{'}355~J per request, reranking adds 106~J and hybrid adds 6~J. HyDE's
marginal cost exceeds the total energy of C0 by a factor of 4.3. The three
modules differ by two orders of magnitude in energy cost.

\begin{figure}[h!]
    \centering
    \includegraphics[width=0.85\textwidth]{feature_attribution_energy.png}
    \caption{Marginal energy cost of each module (delta from C0).}
    \label{fig:feature-attribution}
    % Source: hotpot/2-performance/output/plots/feature_attribution_energy.png
\end{figure}

% --- Energy per token ---

Normalizing by total token count yields energy per token (J/tok), which
controls for additional tokens and enables fairer comparison across setups.
C0 consumes 0.39~J/tok while C7 consumes 1.44~J/tok---a 3.7$\times$ ratio,
compressed from the 5.6$\times$ total energy ratio. HyDE generates
additional tokens and remains the most expensive after normalization, 
confirming that the overhead is structural
(an additional LLM call) rather than a token count artifact.

\begin{figure}[h!]
    \centering
    \begin{subfigure}[t]{0.48\textwidth}
        \centering
        \includegraphics[width=\textwidth]{mean_energy_per_token.png}
        \caption{Mean energy per token}
        \label{fig:energy-per-token-bar}
        % Source: hotpot/2-performance/output/plots/mean_energy_per_token.png
    \end{subfigure}%
    \hfill
    \begin{subfigure}[t]{0.48\textwidth}
        \centering
        \includegraphics[width=\textwidth]{node_energy_per_token_decomposition.png}
        \caption{Energy per token by infrastructure node}
        \label{fig:node-decomposition}
        % Source: hotpot/2-performance/output/plots/node_energy_per_token_decomposition.png
    \end{subfigure}
    \caption{Energy per token by configuration}
    \label{fig:energy-per-token-combined}
\end{figure}

% --- Node-level decomposition prose ---

Figure~\ref{fig:node-decomposition} decomposes energy per token into the three
infrastructure nodes. The LLM node dominates across all configurations; HyDE
(C1) amplifies LLM node cost through an additional generation pass. The
embedding node contributes a smaller fraction and the Milvus node is negligible.\\


HyDE amplifies LLM node energy because it adds a full generation pass before
retrieval. In C1 (HyDE only), the LLM node consumes 1{'}491~J compared to
268~J in C0---a 5.6$\times$ increase on a single node. Reranking primarily
loads the embedding node: C2 raises embedding energy from 43~J (C0) to 66~J,
because the cross-encoder model runs on this node.

% --- Node-device CPU/GPU decomposition ---


Decomposing CPU and GPU energy by infrastructure node reveals that
the absolute energy differs by an order of magnitude across nodes. As expected,
GPU energy dominates only the two model nodes with 80--92\% GPU usage. The
Milvus node is only 42--43\% GPU---vector retrieval relies more on CPU than
on GPU. This aligns with Chen~et~al.~\cite{cpucentric2025}, who report that
CPU accounts for 90.6\% of latency but only 44\% of energy in agentic RAG
workloads. The hardware cost structure differs across pipeline stages, even
when the aggregate CPU/GPU split appears constant (79--81\% GPU across all
configurations).

\begin{figure}[h!]
    \centering
    \includegraphics[width=0.95\textwidth]{node_device_faceted.png}
    \caption{CPU/GPU energy split per infrastructure node}
    \label{fig:node-device}
    % Source: hotpot/2-performance/output/plots/node_device_faceted.png
\end{figure}


% --- Distributional view ---

The HotpotQA dataset could in principle produce
different energy profiles across query subtypes.
Figure~\ref{fig:energy-by-querytype} shows that the configuration ranking
holds consistently across both subtypes.
Bridge queries exhibit wider energy distributions than comparison queries,
particularly in HyDE-enabled configurations. 
Paradoxically, they also consume slightly more
energy and latency on average, yet achieve higher retrieval scores (0--1
scale, higher is better).

% --- Section summary ---

In summary, enabling all modules (C7) costs 5.6$\times$ more energy than the
minimal pipeline (C0). The cost is dominated by HyDE's additional LLM
generation pass, which alone adds 1{'}355~J---more than four times the entire
C0 baseline. Module costs are approximately additive and consistent across
query subtypes. The LLM node carries 84--88\% of total energy, though hardware
utilization differs by node: GPU-dominated for LLM and embedding, CPU-heavy
for vector retrieval. 
% Section~\ref{subsec:latency-energy} 

\begin{figure}[h!]
    \centering
    \includegraphics[width=0.9\textwidth]{energy_by_query_type.png}
    \caption{Energy distribution by query subtype}
    \label{fig:energy-by-querytype}
    % Source: hotpot/2-performance/output/plots/energy_by_query_type.png
\end{figure}

\subsection{Energy-Latency Correlation}
\label{subsec:latency-energy}

% --- Latency by configuration ---
We examine the latency dimension and the relationship between energy and execution time.

The figure below reports the TTFT(left) and response generation time (right). 
the mean TTFT ranges from 982~ms (C0) to 6{'}153~ms (C7); a
6.3$\times$ difference. The latency structure mirrors the energy clusters
from Section~\ref{subsec:energy-breakdown} where C0 and C3 share the lowest TTFT
($\sim$982--1{'}032~ms) followed by reranking configurations, then HyDE configurations 
and then combined configurations around
6{'}109--6{'}153~ms.

The answer generation time is approximately constant across
all configurations ($\sim$112~ms). It confirms that the LLM generation phase
does not vary with retrieval enrichment.

\begin{figure}[h!]
    \centering
    \includegraphics[width=0.9\textwidth]{latency_by_config.png}
    \caption{TTFT and total response time per configuration}
    \label{fig:latency-by-config}
    % Source: hotpot/2-performance/output/plots/latency_by_config.png
\end{figure}

Again, HyDE module imposes the largest TTFT penalty because it requires a complete LLM
generation pass before retrieval can begin. The HyDE adds 4{'}293~ms to TTFT,
reranking adds 825~ms and hybrid adds 50~ms. The combined TTFT of C7
(6{'}153~ms) closely matches the sum of individual deltas
(C0 + HyDE + reranking + hybrid = 982 + 4{'}293 + 825 + 50 = 6{'}150~ms),
confirming that module latencies are again approximately additive.

% --- Cross-dimension cost summary ---

Table~\ref{tab:module-cost} consolidates energy and latency costs per module.
The three modules differ by two orders of magnitude in cost. HyDE
adds 1{'}355~J and 4{'}293~ms; hybrid adds 6~J and 50~ms. This
asymmetry has direct implications for control: hybrid retrieval can be
enabled generously, while HyDE activation must be selective to achieve
meaningful energy savings.

\begin{table}[h!]
\centering
\begin{tabular}{l r r r}
\hline
\textbf{Module} & \textbf{$\Delta$Energy (J)} & \textbf{$\Delta$TTFT (ms)} & \textbf{$\Delta$Total time (ms)} \\
\hline
HyDE              & +1{'}355 & +4{'}293 & +4{'}292 \\
Reranking          & +106     & +825     & +822 \\
Hybrid retrieval   & +6       & +50      & +46 \\
\hline
\end{tabular}
\caption{Per-module marginal cost (C0 delta).}
\label{tab:module-cost}
\end{table}

% --- Energy-latency relationship ---

The parallel structure of energy and latency raises a practical question:
can latency serve as a proxy for energy? Prior work often treats response time
as a stand-in for computational cost; it is straightforward to measure
and directly interpretable. 

To test this assumption, Figure~\ref{fig:energy-latency-regression} plots both total
energy (right panel) and energy per token (left panel) against total response
time for all 9{'}600~sessions. Linear regression of total energy on response
time yields $R^2 = 0.978$ with a slope of 310~J/s---confirming that latency
is a strong first-order proxy. The energy-per-token regression yields a
lower $R^2 = 0.860$, reflecting the additional variance introduced by
token-generation efficiency across configurations.

\begin{figure}[h!]
    \centering
    \includegraphics[width=0.95\textwidth]{energy_latency_regression.png}
    \caption{Energy versus total response time}
    \label{fig:energy-latency-regression}
    % Source: hotpot/2-performance/output/plots/energy_latency_regression.png
    %    across 9{'}600~sessions. Left:energy per token versus: total response time ($R^2 = 0.860$). Right: totalenergy versus total response time ($R^2 = 0.978$). The global correlationis strong, but per-configuration slopes (colored lines) diverge.
\end{figure}

The 12-point gap between the two $R^2$ values (0.978 vs 0.860) shows that
latency is a good but imperfect proxy for per-token energy cost. Configurations that
activate reranking draw less instantaneous power than those that invoke HyDE; a 
cross-encoder pass is less GPU-intensive than a full LLM generation
step. 
%Two sessions with identical latency can therefore differ in energy per
%token.

Most studies report only latency; fewer measure total energy, as power
instrumentation is not straightforward. Rago et al.~\cite{rago2025}
recommended energy per token (J/tok) as an energy metric but listed
it as future work only.
% Our setup implements this metric; the $R^2 = 0.860$
% result confirms that it captures variance that latency alone misses.

% --- Energy-latency residuals ---
\begin{figure}[h!]
    \centering
    \includegraphics[width=0.8\textwidth]{energy_latency_residuals.png}
    \caption{Energy--latency residuals by configuration}
    \label{fig:energy-latency-residuals}
    % Source: hotpot/2-performance/output/plots/energy_latency_residuals.png
\end{figure}

The figure below makes the deviation explicit.
Reranking configurations (C2, C6) consume 81--91~J \emph{less} energy than
predicted by their latency, while non-reranking configurations (C0, C1, C3,
C5) consume 60--90~J \emph{more} than predicted. 

Positive values indicate configurations that consume more energy than predicted by their
latency; negative values indicate power-frugal configurations. The residuals expose two
distinct hardware utilization regimes:

\begin{itemize}
    \item reranking is ``power-frugal'' (lower GPU intensity from the cross-encoder) 
    \item HyDE is ``GPU-intensive'' (sustained GPU utilization for an additional LLM call). 
\end{itemize}

% Power draw (defined as $E_\text{total} / t_\text{total}$) ranges
% from 217~W (C6) to 309~W (C1), a 40\% spread.


Latency is therefore a good proxy for fast energy
approximation or when direct measurement is impractical.
Guo~et~al.~\cite{keplerenergy2026} reach a similar conclusion in their RAG
experiment and caution that ``the relationship is not perfectly linear.'' Our
results do not improve on this approximation, but they quantify the deviation:
two configurations with equal latency can, in our case differ by up to 
40\% in energy. For accurate per-query cost attribution, energy per token is a more precise
metric; accounts for both power draw variation and output length.

% --- Energy per token vs latency ---
\begin{figure}[h!]
    \centering
    \includegraphics[width=0.8\textwidth]{energy_per_token_vs_latency.png}
    \caption{Energy per token versus total response time}
    \label{fig:energy-per-token}
    % Source: hotpot/2-performance/output/plots/energy_per_token_vs_latency.png
    % HyDE-containing configurations cluster in the high-cost region even after normalization.m n
\end{figure} 

Figure~\ref{fig:energy-per-token} illustrates this point. The baseline consumes
0.39~J/token while the full pipeline consumes 1.44~J/token---a 3.7$\times$ ratio.
HyDE configurations remain the most expensive after normalization, confirming 
structural overhead rather than a token count artifact.

% HyDE configurations generate additional tokens (the hypothetical document),
% which compresses the ratio from 5.6$\times$ (total energy) to 3.7$\times$
% (J/tok). 

% --- Section summary ---

Module costs are highly asymmetric: HyDE dominates both energy and latency,
reranking imposes moderate cost and hybrid retrieval is nearly free. Costs are
approximately additive across modules. The energy-latency relationship is
strong ($R^2 = 0.978$ globally) but configuration-dependent: power draw ranges
from 217~W (reranking) to 309~W (HyDE), so latency alone does not fully
determine energy cost. The energy ordering is more robust under per-token
normalization. Next, we ask whether this cost buys
proportional quality improvement.


\subsection{Quality Trade-offs}
\label{subsec:energy-quality}

% --- Block A: Introducing quality ---

We just established that optional modules impose large and asymmetric energy and
latency overhead. A natural intuition is that these modules improve answer
quality and that the cost is therefore justified. To evaluate this claim, we
examine quality outcomes per configuration using two RAGAS metrics:

\begin{itemize}
    \item \textbf{Faithfulness} fraction of claims supported by retrieved context
    \item \textbf{Factual correctness} semantic agreement with the reference answer
\end{itemize}

These two metrics capture complementary failure modes:
hallucination versus factual error.

% --- Block B: Quality scores across configurations ---

Factual correctness ranges from 0.590 (C5) to 0.633 (C6) across all eight
configurations; a spread of 0.043~points. Faithfulness is wider but still
modest, ranging from 0.771 (C1) to 0.830 (C7); a spread of 0.060~points.
Contrast this with the energy dimension: a 5.6$\times$ energy range produces
less than 6\% of quality variation on either metric. 

\begin{figure}[h!]
    \centering
    \includegraphics[width=1.0\textwidth]{quality_heatmap_combined.png}
    \caption{Quality deltas from baseline (C0)}
    \label{fig:feature-attribution-quality}
    % Source: hotpot/3-quality/output/plots/feature_attribution_quality.png
    % per pipeline enhancement across four metrics. Reranking consistently improves quality; HyDE consistently degrades it.
\end{figure}

Within-configuration variance dominates between-configuration differences:
the 95\% confidence interval for factual correctness within any single
configuration ($\pm$0.023) covers more than half the cross-configuration
range (0.043). This indicates that query-level difficulty drives quality more
than module activation does.

% --- Block C: Quality varies by query subtype ---

The Figure~\ref{fig:feature-attribution-quality} (right) splits
factual correctness by querysubtype. Comparison queries score consistently higher than bridge queries
across all configurations with a small gap of 0.08--0.15~points. Within each
subtype, reranking configurations (C2, C4, C6, C7) outperform non-reranking
ones. For bridge queries, C6 (hybrid + reranking) achieves the highest score
(0.588), while C5 (HyDE + hybrid) produces the lowest (0.521). For comparison
queries, C2 (reranking only) leads (0.687), while C1 (HyDE only) scores lowest
(0.652).

HyDE degrades quality on both subtypes: comparing each HyDE configuration with
its non-HyDE equivalent yields
factual correctness deltas between $-$0.001 and $-$0.014~points. The effect is
small but consistently negative, meaning HyDE adds substantial energy cost
while providing no quality return on the workload.

The figure~\ref{fig:feature-attribution-quality} generalizes this per-feature
analysis across four quality metrics. Reranking improves all four metrics
(+0.014~fc, +0.025~faithfulness, +0.017~support~F1, +0.032~context recall).
HyDE degrades all four ($-$0.013~fc, $-$0.021~faithfulness,
$-$0.005~support~F1, $-$0.033~context recall). Hybrid retrieval
shows mixed effects: negative on factual correctness ($-$0.013), support~F1 ($-$0.009)
and context recall ($-$0.031) but near-zero on faithfulness.

% --- Block D: Energy-quality tension quantified ---

Table~\ref{tab:energy-quality} compares energy multipliers with quality deltas
for each module. HyDE (C1 versus C0) increases energy by 5.3$\times$ while
\emph{decreasing} average factual correctness by 0.013~points and faithfulness
by 0.021~points. Reranking (C2 versus C0) increases energy by 1.3$\times$ for
quality gains of 0.014 (fc) and 0.025 (faithfulness). Hybrid retrieval (C3
versus C0) increases energy by 1.02$\times$ with a small negative effect on
factual correctness ($-$0.013) and negligible effect on faithfulness (+0.001).

\begin{table}[h!]
\centering
\begin{tabular}{l r r r}
\hline
\textbf{Module} & \textbf{Energy ratio} & \textbf{$\Delta$ Factual corr.} & \textbf{$\Delta$ Faithfulness} \\
\hline
HyDE (C1)       & 5.3$\times$ & $-$0.013 & $-$0.021 \\
Reranking (C2)   & 1.3$\times$ & $+$0.014 & $+$0.025 \\
Hybrid (C3)      & 1.02$\times$ & $-$0.013 & $+$0.001 \\
All (C7)         & 5.6$\times$ & $+$0.016 & $+$0.038 \\
\hline
\end{tabular}
\caption{Energy--quality trade-off per module (C0 delta).}
\label{tab:energy-quality}
\end{table}

The full pipeline (C7) produces a quality improvement of 0.016~points in
factual correctness and 0.038~points in faithfulness at a cost of
5.6$\times$ energy. The quality gain is driven entirely by reranking; HyDE
contributes no positive effect. The relationship between energy investment and
quality return is nonlinear: the first 2\% energy increase (hybrid) yields no
quality improvement, while an additional 31\% (reranking) produces +0.014~fc
and +0.025~faithfulness. The remaining 330\% energy increase from HyDE adds no
further quality benefit.

Figure~\ref{fig:quality-vs-energy-scatter} visualizes this tension. The
x-axis reports mean energy per token (J/tok) rather than total energy to
control for token-count differences introduced by HyDE's additional
generation pass (Section~\ref{subsec:latency-energy}). Two clusters emerge:
non-HyDE configurations (C0, C2, C3, C6) occupy the low-cost region
(0.39--0.51~J/tok), while HyDE configurations (C1, C4, C5, C7) cluster
above 1.36~J/tok with no systematic quality advantage. The Pareto frontier
contains C6 (hybrid + reranking, 0.633~fc at 0.51~J/tok) and C2 (reranking
only, 0.630~fc at 0.50~J/tok).

\begin{figure}[h!]
    \centering
    \includegraphics[width=0.85\textwidth]{quality_vs_energy.png}
    \caption{Factual correctness versus energy per token}
    \label{fig:quality-vs-energy-scatter}
    % Source: hotpot/3-quality/output/plots/quality_vs_energy.png
\end{figure}

% --- Block E: Section summary ---

Quality differences across configurations are small (less than 0.06~points on
both metrics) and vary by subtype. Energy differences are large
(5.6$\times$). HyDE produces the worst cost-to-quality ratio: it increases
energy by 5.3$\times$ while degrading both quality metrics. Reranking offers
the best ratio: moderate cost (1.3$\times$) for consistent quality improvement.

In the context of the HotpotQA dataset, HyDE should be disabled by default.
Across all analyses, HyDE is unambiguously negative for this workload:
5.3$\times$ energy, +4{'}293\,ms latency, $-$0.013~factual correctness
and $-$0.021~faithfulness. The effect is consistent across query subtypes
and all four quality metrics. The immediate practical recommendation is to
disable HyDE and operate in the \{C0, C2, C3, C6\}~space.

Within this reduced space, C6 (hybrid + reranking) sits on the Pareto
frontier at 0.633~factual correctness and 435\,J per request---1.37$\times$
C0's energy for the highest quality among non-HyDE configurations. C2
(reranking only) achieves 0.630~fc at 424\,J. The choice between
always-C6 and always-C0 is the most actionable deployment decision: it
requires zero routing infrastructure and captures the quality--energy
trade-off at its sharpest point.

Section~\ref{subsec:hypothesis-validation} synthesizes all three dimensions.


\subsection{Hypothesis Validation}
\label{subsec:hypothesis-validation}

% --- Module-level summary ---

The three preceding sections converge on one pattern: cost is unconditional
but benefit is conditional. Each module exhibits a distinct profile:

\begin{itemize}
    \item \textbf{HyDE}: highest cost (5.3$\times$ energy, 4{'}293~ms
          latency), negative quality effect ($-$0.013~fc, $-$0.021~faith).
    \item \textbf{Reranking}: moderate cost (1.3$\times$ energy, 825~ms
          latency), consistent quality gain (+0.014~fc, +0.025~faith).
    \item \textbf{Hybrid retrieval}: negligible cost (1.02$\times$ energy,
          50~ms latency), near-zero quality effect.
\end{itemize}

% --- The "trivial vs non-trivial" pivot ---

If all modules degraded quality or left it unchanged, the solution would be
straightforward: disable them and collect the energy savings. HyDE does so;
it costs more energy and returns no
quality benefit. Reranking does not; it consistently improves both
faithfulness and factual correctness despite a modest magnitude.
Also, no single configuration dominates across all queries and subtypes.

However, it could be that a module that helps \emph{on average} does not necessarily help every
query and a module that hurts on average may still help specific ones. Whether
per-query patterns exist and whether they are predictable from the query
alone is yet unknown. If module benefit varies across
queries but follows no learnable pattern, then no routing model can outperform
a static rule. The next chapter formalizes this as a decision
problem and investigates whether the available signals support per-query
configuration selection.
